Here is the motivation letter:

\documentclass[11pt]{letter}

\begin{document}

\opening{\noindent June 30, 2026}
Dear Dr. Fanny Dufosse,

I am excited to apply for the PhD position "The levers of sobriety and conviviality and their influence on algorithmic decision models" at the University of Toulouse under the supervision of Inria.

As a highly motivated AI & Data Science Engineer, I am thrilled about the opportunity to contribute to the SOCLOUD project, which aims to design sustainable and open data centers that respect planetary boundaries and are socially equitable. With my Master's degree in Computer Science and strong programming skills, I am confident that my background and experience make me an ideal candidate for this position.

During my previous project, "Cardiovascular Medical Digital Twin", I had the opportunity to work on several innovative achievements that demonstrate my ability to develop robust and efficient solutions. Specifically, I would like to highlight two of these achievements that map perfectly to the requirements of this PhD offer.

Firstly, my implementation of the Missingness-Aware Gate allowed me to dynamically handle sensor dropouts and missing data in real-time, which is crucial for designing sustainable data centers that can handle uncertainty and variability. This achievement demonstrates my ability to develop robust algorithms that can effectively deal with imperfect data, a key aspect of this PhD project.

Secondly, my development of the Medical_MambaTS linear-time state space model showcases my expertise in developing efficient and low-latency solutions. As the SOCLOUD project aims to explore the possibilities of designing sustainable and open data centers that respect planetary boundaries, I believe that my experience with MambaTS can be leveraged to develop novel algorithmic decision models that take into account energy management.

I am excited about the prospect of joining the University of Toulouse and contributing to the SOCLOUD project. Thank you for considering my application. I look forward to the opportunity to discuss this further.

Sincerely,

Ismail Aissaoui

\closing{\noindent}

\end{document}